\bigskip
\noindent Autres pistes de l'onglet Findings pour ce thème, dans des cotes qui n'ont pas de section ici :\par

\begin{encadre}
\begin{theoreme*}[Grothendieck, pages 17 à 23, factorisation radicielle-étale\piste{126-radicial-etale-factorisation}]\nonverifie
Soit $f : X \to Y$ fini et universellement ouvert, $Y$ localement noethérien, dont toutes les fibres ont le même nombre de points géométriques. Alors $f$ se factorise en un morphisme radiciel fini surjectif $X \to X^{\text{ét}}$ suivi d'un revêtement étale fini $X^{\text{ét}} \to Y$, universel parmi les $Y$-morphismes de $X$ vers les $Y$-schémas étales.
\end{theoreme*}
\end{encadre}
{\small Les feuillets en esquissent une démonstration incomplète : la section ouverte et fermée de la page~19 est marquée \q{[vérif]}, la prose des lemmes des pages 18, 20, 21 et~22 est illisible, et la page~23 s'arrête sur \q{À formuler}.\par}

\begin{encadre}
\begin{theoreme*}[D'après les pages 4 à 7, forme complétée par l'édition, quotient d'une déformation nilpotente\piste{126-nilpotent-deformation-quotient}]\nonverifie
Soit $R \rightrightarrows X$ une relation d'équivalence plate dont la réduction $R_0 \rightrightarrows X_0$ modulo un idéal nilpotent est effective, de quotient fpqc quasi-séparé $X_0 \to Y_0$. Alors $R$ est effective, son quotient $Y$ est un épaississement nilpotent de $Y_0$ et $X \to Y$ est plat.
\end{theoreme*}
\end{encadre}
{\small La page porte la conclusion et la méthode, le relèvement par l'annulation de $\check H^0(X_0/Y_0, \mathcal H^1)$ ; les hypothèses, en grande partie illisibles, sont reconstituées par l'édition, de même que la réduction au cas $\mathcal J^2 = 0$ et la platitude.\par}

\begin{encadre}
\begin{theoreme*}[Grothendieck, pages 26 à 29, descente générique\piste{126-generic-descent}]\nonverifie
Soient $S$ noethérien, $T \to S$ plat de type fini et $F$ une catégorie fibrée sur les $S$-schémas de type fini pour laquelle les morphismes plats surjectifs sont de $F$-descente (pleine fidélité seulement). Un isomorphisme $\varphi : p_1^*X \simeq p_2^*X$ défini sur un ouvert dense de $T \times_S T$ et vérifiant la condition de cocycle sur un ouvert dense de $T \times_S T \times_S T$ se prolonge, après restriction de $T$ à un ouvert dense, en une donnée de descente qui coïncide avec $\varphi$ là où les deux sont définies.
\end{theoreme*}
\end{encadre}
\begin{demopage}[pages 26 à 29]
On réduit $\mathcal U$ et $V$ à des ouverts, et $T$ à un ouvert dense au-dessus duquel les fibres de $\mathcal U$ et de $V$ sont denses (p.~26). Sur $W = p_{12}^{-1}(\mathcal U) \cap p_{23}^{-1}(\mathcal U)$, la page définit $\psi : \pi^*X_1 \simeq \pi^*X_2$, $\pi = p_{13}$, en composant $\varphi$ par un point médian (p.~27). Sa compatibilité avec $\pi$ se lit sur une partie $W'$ de $T^5$ par le cocycle ; elle demande que $W' \to W \times_{T \times_S T} W$ soit surjectif, et la page l'admet \q{sauf erreur} (p.~28). $W \to T \times_S T$ étant plat et surjectif, $\psi$ descend en $\varphi'$, dont le cocycle se vérifie sur une partie $W'''$ de $T^6$ (pp.~28--29), et $\varphi' = \varphi$ sur $\mathcal U$ car $V \to \mathcal U$ est surjectif (p.~29). La lecture réordonne les restrictions.
\end{demopage}
